%% ============================================================
%% LECCIÓN 1.2: LÍMITES, CONTINUIDAD Y DERIVADAS PARCIALES
%% Cálculo III — UIS (20254)
%% Contenido basado en LPS_L1.2.pdf y Constitución de lecciones
%% Estilo: Apostol unificado, con entornos LaTeX personalizados
%% ============================================================

\chapter{Límites y continuidad en $\mathbb{R}^n$}

La teoría de funciones de una variable nos enseñó que el comportamiento
de una función cerca de un punto se captura mediante el límite, y que
la regularidad de ese comportamiento se expresa con la continuidad.
También aprendimos que la razón de cambio instantánea de una función
es la derivada. En esta lección extenderemos estas tres nociones
fundamentales —límite, continuidad y derivada— a funciones de dos o
más variables. La extensión no es meramente formal: el hecho de que
el dominio sea ahora un espacio de dimensión superior introduce una
novedad cualitativa: el punto de aproximación puede alcanzarse a lo
largo de infinitas trayectorias distintas. En una variable solo hay
dos direcciones (izquierda y derecha); en dos variables hay infinitas
curvas que convergen al mismo punto, y la existencia del límite exige
que el valor límite sea el mismo a lo largo de \emph{todas} ellas.

\begin{rem}
El hilo conductor de esta lección es el siguiente principio: para
funciones de varias variables, todo concepto que dependa de la
``proximidad'' se formula en términos de la \textbf{distancia
euclidiana} entre puntos. El límite exige que la función se acerque a
un valor $L$ cuando la distancia entre $(x,y)$ y $(a,b)$ tiende a cero,
\emph{independientemente de la dirección desde la que se aproxime}.
Esta independencia direccional es la novedad central y la principal
fuente de dificultad.
\end{rem}

% ============================================================
% SECCIÓN 1.3: LÍMITE Y CONTINUIDAD DE UNA FUNCIÓN DE VARIAS VARIABLES
% ============================================================

\section{Límite y continuidad de una función de varias variables}

Empezamos con el concepto de límite para funciones de dos variables,
que luego extenderemos a tres o más.

\begin{definition}[Límite de una función de dos variables]
Sea $f$ una función definida en un disco abierto centrado en $(a,b)$,
excepto posiblemente en $(a,b)$. Decimos que
\[
\lim_{(x,y)\to(a,b)} f(x,y) = L
\]
si para todo $\varepsilon>0$ existe $\delta>0$ tal que
\[
0 < \sqrt{(x-a)^2 + (y-b)^2} < \delta
\quad\Longrightarrow\quad
|f(x,y)-L| < \varepsilon.
\]
\end{definition}

Esta definición es la análoga bidimensional de la definición
$\varepsilon$-$\delta$ para una variable: la distancia entre $(x,y)$ y
$(a,b)$ se mide ahora con la norma euclidiana en $\mathbb{R}^2$, y
la condición $0<\sqrt{\cdots}$ excluye al punto $(a,b)$ del análisis,
pues el límite no depende del valor de $f$ en ese punto.

\begin{rem}
La definición exige que la función esté definida en \emph{algún}
disco perforado alrededor de $(a,b)$; no necesita estar definida
exactamente en $(a,b)$. Esta es la misma convención que en una
variable.
\end{rem}

\begin{rem}
Para funciones de tres variables, la definición es idéntica
sustituyendo la distancia en $\mathbb{R}^3$:
\[
\lim_{(x,y,z)\to(a,b,c)} f(x,y,z) = L
\]
si para todo $\varepsilon>0$ existe $\delta>0$ tal que
$0<\sqrt{(x-a)^2+(y-b)^2+(z-c)^2}<\delta$ implica
$|f(x,y,z)-L|<\varepsilon$.
\end{rem}

La diferencia fundamental con el caso unidimensional es que ahora el
punto $(a,b)$ puede ser alcanzado por \emph{infinitas trayectorias}
distintas en el plano. En una variable, solo hay dos caminos de
aproximación: por la izquierda o por la derecha. En dos variables,
hay infinitos: rectas de cualquier pendiente, curvas como parábolas,
cúbicas, etc., que pasen por $(a,b)$. Para que el límite exista, el
valor de $f$ debe tender al mismo $L$ a lo largo de \emph{todas}
ellas. Si encontramos dos trayectorias que producen valores límite
distintos, el límite no existe.

\begin{rem}
Este principio es análogo al de los límites laterales en una variable:
$\lim_{x\to a}f(x)$ existe si y solo si $\lim_{x\to a^-}f(x)$ y
$\lim_{x\to a^+}f(x)$ existen y coinciden. En varias variables, los
límites laterales se reemplazan por límites a lo largo de trayectorias,
y hay infinitas trayectorias posibles.
\end{rem}

\begin{pasos}[Para determinar si un límite de dos variables existe]
\begin{enumerate}
    \item \textbf{Prueba de trayectorias.} Si la función es sencilla,
    intentar evaluar a lo largo de rectas $y=mx$ (o curvas $y=x^2$,
    $x=y^2$, etc.) que pasen por el punto. Si dos trayectorias
    conducen a valores distintos, el límite no existe.
    \item \textbf{Sustitución directa.} Si la función es continua en el
    punto (por ejemplo, polinómica), el límite es $f(a,b)$.
    \item \textbf{Manipulación algebraica.} Factorizar, racionalizar o
    simplificar para eliminar indeterminaciones.
    \item \textbf{Coordenadas polares.} Si el punto de aproximación es
    el origen, usar $x=r\cos\theta$, $y=r\sin\theta$. Si el límite
    cuando $r\to 0$ es independiente de $\theta$, el límite existe; si
    depende de $\theta$, no existe.
    \item \textbf{Teorema del emparedado.} Acotar $|f(x,y)-L|$ por
    una función que tiende a cero.
\end{enumerate}
\end{pasos}

\begin{example}[Límite que no existe por trayectorias distintas]
Sea
\[
f(x,y) = \frac{x^2 y}{x^4 + y^2}.
\]
Estudiar $\lim_{(x,y)\to(0,0)} f(x,y)$.
\begin{myproof}
\textbf{Paso 1. Prueba de trayectorias.} Evaluamos a lo largo de dos
trayectorias distintas que pasan por el origen.

\textit{Trayectoria 1:} $y=x$. Entonces
\[
f(x,x) = \frac{x^2\cdot x}{x^4+x^2} = \frac{x^3}{x^2(x^2+1)}
        = \frac{x}{x^2+1}.
\]
Cuando $x\to 0$,
\[
\lim_{x\to 0} \frac{x}{x^2+1} = 0.
\]

\textit{Trayectoria 2:} $y=x^2$. Entonces
\[
f(x,x^2) = \frac{x^2\cdot x^2}{x^4+(x^2)^2}
          = \frac{x^4}{x^4+x^4} = \frac{1}{2}.
\]
Cuando $x\to 0$,
\[
\lim_{x\to 0} \frac{1}{2} = \frac{1}{2}.
\]

\textbf{Conclusión.} Como las dos trayectorias producen límites
distintos ($0$ y $\frac12$), el límite no existe en $(0,0)$.
$\boxed{\text{No existe}}$
\end{myproof}
\end{example}

\begin{example}[Límite que no existe por trayectorias paramétricas]
Sea
\[
f(x,y) = \frac{x^3 y}{x^6 + y^2}.
\]
Estudiar $\lim_{(x,y)\to(0,0)} f(x,y)$.
\begin{myproof}
\textbf{Paso 1. Prueba de trayectorias.} Evaluamos a lo largo de dos
trayectorias.

\textit{Trayectoria 1:} $y=x^3$. Entonces
\[
f(x,x^3) = \frac{x^3\cdot x^3}{x^6+(x^3)^2}
          = \frac{x^6}{x^6+x^6} = \frac{1}{2}.
\]
Así, $\lim_{x\to 0} f(x,x^3)=\frac12$.

\textit{Trayectoria 2:} $y=x$. Entonces
\[
f(x,x) = \frac{x^3\cdot x}{x^6+x^2}
        = \frac{x^4}{x^2(x^4+1)} = \frac{x^2}{x^4+1}.
\]
Cuando $x\to 0$,
\[
\lim_{x\to 0} \frac{x^2}{x^4+1} = 0.
\]

\textbf{Conclusión.} Las trayectorias $y=x^3$ e $y=x$ dan límites
distintos ($\frac12$ y $0$), por lo que el límite no existe.
$\boxed{\text{No existe}}$
\end{myproof}
\end{example}

\begin{example}[Límite que existe mediante coordenadas polares]
Sea
\[
f(x,y) = \frac{x^2 y}{x^2 + y^2}.
\]
Estudiar $\lim_{(x,y)\to(0,0)} f(x,y)$.
\begin{myproof}
\textbf{Paso 1. Coordenadas polares.} Hacemos $x=r\cos\theta$,
$y=r\sin\theta$. Entonces
\[
f(r\cos\theta, r\sin\theta)
= \frac{(r\cos\theta)^2(r\sin\theta)}
       {(r\cos\theta)^2+(r\sin\theta)^2}
= \frac{r^3\cos^2\theta\sin\theta}{r^2}
= r\cos^2\theta\sin\theta.
\]

\textbf{Paso 2. Límite cuando $r\to 0$.}
Para cualquier $\theta$ fijo,
\[
\lim_{r\to 0} r\cos^2\theta\sin\theta = 0,
\]
pues $|\cos^2\theta\sin\theta| \leq 1$ y $r\to 0$.

\textbf{Paso 3. Conclusión.}
El límite es $0$, independiente de $\theta$. Por tanto,
$\displaystyle \lim_{(x,y)\to(0,0)} f(x,y) = 0$.
$\boxed{0}$
\end{myproof}
\end{example}

\begin{rem}
El método de coordenadas polares es particularmente útil cuando la
función contiene expresiones $x^2+y^2$ en el denominador, pues esa
expresión se convierte en $r^2$, simplificando el álgebra. Sin
embargo, hay que tener cuidado: que el límite cuando $r\to 0$ no
dependa de $\theta$ \emph{para cada $\theta$ fijo} no siempre
garantiza la existencia del límite; la definición $\varepsilon$-$\delta$
exige que la convergencia sea uniforme en $\theta$. Para funciones
continuas en $\theta$ y acotadas, como en el ejemplo anterior, la
convergencia uniforme se sigue de manera inmediata. En casos más
sutiles, puede ser necesario un análisis adicional.
\end{rem}

\begin{example}[Límite que no existe pese a ser constante en rectas]
Sea
\[
f(x,y) = \frac{x^2 y^2}{x^4 + y^4}.
\]
Evaluar a lo largo de rectas $y=mx$:
\[
f(x,mx) = \frac{x^2(mx)^2}{x^4+(mx)^4}
        = \frac{m^2 x^4}{x^4(1+m^4)}
        = \frac{m^2}{1+m^4},
\]
que depende de $m$. Por tanto, el límite no existe.
\end{example}

\begin{rem}
Este ejemplo muestra que la prueba de trayectorias puede ser más
sensible que las rectas. En la función anterior, todas las rectas dan
límites que dependen de la pendiente, lo que es suficiente para
concluir la no existencia. Sin embargo, hay funciones donde todas las
rectas dan el mismo límite, pero curvas no rectas dan otro; en ese
caso, la prueba de rectas no es suficiente y hay que explorar
trayectorias curvas.
\end{rem}

% ------------------------------------------------------------
% PROPIEDADES DE LOS LÍMITES
% ------------------------------------------------------------

Las propiedades algebraicas de los límites se extienden naturalmente a
varias variables: si $\lim f=L$ y $\lim g=K$, entonces
\[
\begin{aligned}
\lim (f+g) &= L+K,\\
\lim (f-g) &= L-K,\\
\lim (f\cdot g) &= L\cdot K,\\
\lim \frac{f}{g} &= \frac{L}{K} \quad (K\neq 0),\\
\lim (k f) &= kL \quad (k\in\mathbb{R}).
\end{aligned}
\]
Estas propiedades permiten evaluar límites de funciones que son
combinaciones de funciones elementales mediante sustitución directa,
siempre que la función esté definida en el punto y no haya
indeterminación.

\begin{example}[Límite por sustitución directa]
Calcular
\[
\lim_{(x,y)\to(\pi,1)} \frac{\cos(xy)}{y^2+1}.
\]
\begin{myproof}
\textbf{Paso 1:} La función $\dfrac{\cos(xy)}{y^2+1}$ es cociente de funciones
continuas en $(\pi,1)$, y el denominador no se anula: $1^2+1=2\neq 0$.
Por tanto, podemos evaluar por sustitución directa.

\textbf{Paso 2:} Sustituimos $(x,y)=(\pi,1)$:
\[
\lim_{(x,y)\to(\pi,1)} \frac{\cos(xy)}{y^2+1}
= \frac{\cos(\pi\cdot 1)}{1^2+1}
= \frac{-1}{2}.
\]

El límite es $\displaystyle\boxed{-\frac{1}{2}}$.
\end{myproof}
\end{example}

% ------------------------------------------------------------
% TEOREMA DEL EMPAREDADO (SÁNDWICH) EN VARIAS VARIABLES
% ------------------------------------------------------------

\begin{theorem}[Teorema del emparedado para funciones de dos variables]
Si existe un disco perforado centrado en $(a,b)$ donde
\[
|f(x,y)-L| \leq g(x,y)
\]
y además $\lim_{(x,y)\to(a,b)} g(x,y) = 0$, entonces
\[
\lim_{(x,y)\to(a,b)} f(x,y) = L.
\]
\end{theorem}

\begin{proof}
Dado $\varepsilon>0$, como $g\to 0$, existe $\delta>0$ tal que
$0<\sqrt{(x-a)^2+(y-b)^2}<\delta$ implica $g(x,y)<\varepsilon$.
Entonces $|f(x,y)-L|\leq g(x,y)<\varepsilon$. Luego el límite existe
y es $L$.
\end{proof}

\begin{example}[Uso del teorema del emparedado]
Calcular
\[
\lim_{(x,y)\to(0,0)} \frac{x^2 y^2}{x^2+y^2}.
\]
\begin{myproof}
\textbf{Paso 1:} Establecemos una cota superior. Dado que $x^2 y^2 \leq (x^2+y^2)^2$
(pues $(x^2+y^2)^2 = x^4+2x^2y^2+y^4 \geq x^2y^2$ para $x^2,y^2\geq 0$), se tiene:
\[
0 \leq \frac{x^2 y^2}{x^2+y^2} \leq \frac{(x^2+y^2)^2}{x^2+y^2} = x^2+y^2.
\]

\textbf{Paso 2:} Aplicamos el teorema del emparedado. Como
$\lim_{(x,y)\to(0,0)}(x^2+y^2)=0$, concluimos
\[
\lim_{(x,y)\to(0,0)} \frac{x^2 y^2}{x^2+y^2} = \boxed{0}.
\]
\end{myproof}
\end{example}

% ------------------------------------------------------------
% CONTINUIDAD
% ------------------------------------------------------------

\begin{definition}[Continuidad en un punto]
Una función $f$ de dos variables es \textbf{continua en $(a,b)$} si:
\begin{enumerate}
    \item $f(a,b)$ está definida,
    \item $\lim_{(x,y)\to(a,b)} f(x,y)$ existe,
    \item $\lim_{(x,y)\to(a,b)} f(x,y) = f(a,b)$.
\end{enumerate}
Si alguna de estas condiciones falla, $f$ tiene una \textbf{discontinuidad}
en $(a,b)$.
\end{definition}

La continuidad en varias variables conserva la idea intuitiva de que
la gráfica no presenta saltos, huecos o interrupciones. En dos
variables, la gráfica es una superficie; la continuidad en un punto
significa que la superficie no se rompe en ese punto.

\begin{rem}
Como en una variable, todas las funciones polinómicas, racionales
(con denominador no nulo), exponenciales, logarítmicas, trigonométricas
y sus inversas son continuas en todo su dominio.
\end{rem}

\begin{definition}[Continuidad en conjuntos]
\begin{itemize}
    \item $f$ es continua en un \textbf{conjunto abierto} $R\subset\mathbb{R}^2$
    si es continua en todo punto de $R$.
    \item $f$ es continua en un \textbf{conjunto cerrado} $R$ si es
    continua en todo punto interior y, en los puntos frontera, el
    límite desde dentro de la región coincide con el valor de la
    función.
\end{itemize}
\end{definition}

\begin{rem}
Las definiciones de conjunto abierto, cerrado, punto interior y punto
frontera en $\mathbb{R}^2$ son las mismas que en $\mathbb{R}$, pero
con discos (circunferencias) en lugar de intervalos. Un disco abierto
es el interior de una circunferencia; un disco cerrado incluye su
borde. Un punto es interior si existe un disco abierto centrado en él
totalmente contenido en el conjunto. Un punto es frontera si todo
disco centrado en él contiene puntos del conjunto y puntos fuera de él.
\end{rem}

\begin{theorem}[Composición de funciones continuas]
Si $f$ es continua en $(a,b)$ y $g$ es continua en $f(a,b)$, entonces
la composición $g(f(x,y))$ es continua en $(a,b)$.
\end{theorem}

\begin{proof}
La demostración es análoga al caso unidimensional: se aplica la
definición de continuidad de $g$ en $f(a,b)$ y luego la de $f$ en
$(a,b)$ para controlar la proximidad.
\end{proof}

\begin{example}[Continuidad de una función definida por partes]
Estudiar la continuidad de
\[
f(x,y)=
\begin{cases}
\dfrac{x^2 y}{x^2+y^2}, & (x,y)\neq (0,0),\\[8pt]
0, & (x,y)=(0,0).
\end{cases}
\]
\begin{myproof}
\textbf{Paso 1:} Para $(x,y)\neq(0,0)$, $f$ es cociente de funciones continuas
con denominador $x^2+y^2>0$, luego es continua en
$\mathbb{R}^2\setminus\{(0,0)\}$.

\textbf{Paso 2:} Verificamos la continuidad en $(0,0)$ con coordenadas polares
($x=r\cos\theta$, $y=r\sin\theta$):
\[
\frac{x^2 y}{x^2+y^2} = \frac{r^3\cos^2\theta\sin\theta}{r^2}
= r\cos^2\theta\sin\theta \;\xrightarrow{r\to 0}\; 0.
\]
Luego $\lim_{(x,y)\to(0,0)} f(x,y)=0=f(0,0)$, así que $f$ es continua en $(0,0)$.

\textbf{Paso 3:} $f$ es continua en todo punto de $\mathbb{R}^2$, por tanto
$\boxed{f \text{ es continua en todo } \mathbb{R}^2}$.
\end{myproof}
\end{example}

% ------------------------------------------------------------
% EXTENSIÓN A TRES VARIABLES
% ------------------------------------------------------------

\begin{definition}[Continuidad para tres variables]
Una función $f$ de tres variables es continua en $(a,b,c)$ si
\[
\lim_{(x,y,z)\to(a,b,c)} f(x,y,z) = f(a,b,c).
\]
\end{definition}

La definición de límite para tres variables es análoga usando la
distancia en $\mathbb{R}^3$: $0<\sqrt{(x-a)^2+(y-b)^2+(z-c)^2}<\delta$.
La continuidad de las funciones elementales se conserva.

\begin{example}[Límite en tres variables]
Calcular $\displaystyle\lim_{(x,y,z)\to(0,0,0)} \frac{x^2 - y^2}{x^2+y^2+z^2}$.
\begin{myproof}
\textbf{Paso 1. Prueba de trayectorias.}

\textit{Trayectoria 1:} $x=y$, $z=0$. Entonces
\[
f(x,x,0) = \frac{x^2 - x^2}{x^2+x^2+0} = 0.
\]

\textit{Trayectoria 2:} $x=0$, $z=0$. Entonces
\[
f(0,y,0) = \frac{0 - y^2}{0+y^2+0} = -1.
\]

\textbf{Conclusión.} Como las trayectorias dan límites distintos,
el límite no existe. $\boxed{\text{No existe}}$
\end{myproof}
\end{example}

\section{Problemas propuestos}

\begin{prob}
Calcule los siguientes límites o demuestre que no existen:
\begin{multicols}{2}
\begin{enumerate}
    \item $\displaystyle\lim_{(x,y)\to(0,0)} \frac{x^2 y}{x^2+y^2}$
    \item $\displaystyle\lim_{(x,y)\to(0,0)} \frac{x^4 y^4}{(x^2+y^4)^3}$
    \item $\displaystyle\lim_{(x,y)\to(0,0)} \frac{x^3+y^3}{x^2+y^2}$
    \item $\displaystyle\lim_{(x,y)\to(0,0)} \frac{x^2 y^2}{x^4+y^4}$
    \item $\displaystyle\lim_{(x,y)\to(0,0)} \frac{x^2 y}{x^4+y^2}$
    \item $\displaystyle\lim_{(x,y)\to(0,0)} \frac{\sin(x^2+y^2)}{x^2+y^2}$
    \item $\displaystyle\lim_{(x,y)\to(0,0)} \frac{1-\cos(xy)}{x^2 y^2}$
    \item $\displaystyle\lim_{(x,y)\to(0,0)} \frac{x y}{\sqrt{x^2+y^2}}$
\end{enumerate}
\end{multicols}
\end{prob}

\begin{prob}
Estudie la continuidad de las siguientes funciones en todo su dominio:
\begin{enumerate}
    \item $f(x,y)=\begin{cases}
    \dfrac{x^2-y^2}{x^2+y^2}, & (x,y)\neq(0,0),\\
    0, & (x,y)=(0,0).
    \end{cases}$
    \item $f(x,y)=\begin{cases}
    \dfrac{x^2 y}{x^2+y^2}, & (x,y)\neq(0,0),\\
    0, & (x,y)=(0,0).
    \end{cases}$
    \item $f(x,y)=\begin{cases}
    \dfrac{x^3 y}{x^6+y^2}, & (x,y)\neq(0,0),\\
    0, & (x,y)=(0,0).
    \end{cases}$
\end{enumerate}
\end{prob}

\begin{prob}
Calcule $\displaystyle\lim_{(x,y)\to(0,0)} \frac{x^2+y^2}{\sqrt{x^2+y^2+1}-1}$
usando coordenadas polares.
\end{prob}

\begin{prob}
Demuestre que $\displaystyle\lim_{(x,y)\to(0,0)} \frac{x^2 y^2}{x^2+y^2}=0$
usando el teorema del emparedado.
\end{prob}

\begin{prob}
Sea $f(x,y)=\dfrac{x^3 y}{x^4+y^2}$ para $(x,y)\neq(0,0)$.
\begin{enumerate}
    \item Demuestre que a lo largo de cualquier recta $y=mx$, el límite
    cuando $(x,y)\to(0,0)$ es $0$.
    \item Demuestre que a lo largo de la parábola $y=x^2$, el límite
    es distinto. Concluya que el límite no existe.
\end{enumerate}
\end{prob}


% ============================================================
% Cap. 19 — Límites y continuidad en Rⁿ
% 16 problemas unicos: 16 listos para integrar ahora, 0 esperan pasada Claude Code (pstricks)
% ============================================================

% >>> LISTOS PARA INTEGRAR AHORA <<<

% [Original: líneas 242-259 | solución completa | sin figura]
\begin{prob}
Sea $f(x,y)=\dfrac{x-y}{x+y}.$
\begin{enumerate}[(a)]
    \item Calcule $\displaystyle\lim_{(x,y)\to (0,0)}f(x)$ tomando la trayectoria $y=kx.$
    \item ¿Qué conclusión puede extraer?
\end{enumerate}
\begin{myproof}
Sea $f(x,y)=\dfrac{x-y}{x+y}.$ Observe que:
\begin{enumerate}[$a)$]
\item Tomando la trayectoria $y=kx$ el límite queda así:
\begin{align*}
\displaystyle\lim_{(x,y)\to (0,0)} \dfrac{x-y}{x+y} {{=}\atop {\begin{array}{c}
\uparrow\\ \text{ si $y=kx$}\end{array}} } \displaystyle\lim_{x\to 0}\dfrac{x-kx}{x+kx}=\displaystyle\lim_{x\to 0}\dfrac{x(1-k)}{x(1+k)}=\dfrac{1-k}{1+k}.
\end{align*}
\item El límite no existe, pues depende del valor de $k\in \mathbb{R}$.
\end{enumerate}
\end{myproof}
\end{prob}

% [Original: líneas 285-308 | solución completa | sin figura]
\begin{prob}
Sea $f(x,y)=\dfrac{xy+y^3}{x^2+y^2}.$
\begin{enumerate}[(a)]
    \item Calcule $\displaystyle\lim_{(x,y)\to (0,0)}f(x)$ tomando la trayectoria $y=x.$
    \item Calcule $\displaystyle\lim_{(x,y)\to (0,0)}f(x)$ tomando la trayectoria $y=0.$
    \item ¿Qué conclusión puede extraer?
\end{enumerate}
\begin{myproof}
Sea $f(x,y)=\dfrac{xy+y^3}{x^2+y^2}.$ Observe que:
\begin{enumerate}[$a)$]
\item Tomando la trayectoria $y=x$ el límite queda así:
\begin{align*}
\displaystyle\lim_{(x,y)\to (0,0)} \dfrac{xy+y^3}{x^2+y^2} {{=}\atop {\begin{array}{c}
\uparrow\\ \text{ si $y=x$}\end{array}} } \displaystyle\lim_{x\to 0}\dfrac{x^2+x^3}{x^2+x^2}=\displaystyle\lim_{x\to 0}\dfrac{x^2(1+x)}{2x^2}=\dfrac{1}{2}\lim_{x\to 0}(1+x)=\dfrac{1}{2}.
\end{align*}
\item Tomando la trayectoria $y=0$ el límite queda así:
\begin{align*}
\displaystyle\lim_{(x,y)\to (0,0)} \dfrac{xy+y^3}{x^2+y^2} {{=}\atop {\begin{array}{c}
\uparrow\\ \text{ si $y=0$}\end{array}} } \displaystyle\lim_{x\to 0}\dfrac{x(0)+0^3}{x^2+0^2}=\displaystyle\lim_{x\to 0}\dfrac{0}{x^2}=0.
\end{align*}
\item El límite no existe, pues dos trayectorias determinan un valor distinto.
\end{enumerate}
\end{myproof}
\end{prob}

% [Original: líneas 333-343 | solución completa | sin figura]
\begin{prob}
Sea $f(x,y)=\dfrac{\arctan(xy)}{xy}.$ Sobre esta función se conoce que $1-\dfrac{x^2y^3}{3}<\dfrac{\arctan (xy)}{xy}<1.$ Calcule $\displaystyle\lim_{(x,y)\to (0,0)}f(x,y).$
\begin{myproof}
Sea $f(x,y)=\dfrac{\arctan(xy)}{xy}.$ Observe que:
\begin{enumerate}[$a)$]
\item $\displaystyle\lim_{(x,y)\to (0,0)} 1-\dfrac{x^2y^3}{3} = 1-\displaystyle\lim_{(x,y)\to (0,0)} \dfrac{x^2y^3}{3}=1-0=1.$
\item $\displaystyle\lim_{(x,y)\to (0,0)} 1=1.$
\item Por medio del teorema del emparedado el límite de la función $f(x,y)=\dfrac{\arctan(xy)}{xy}$  existe y es igual a $1.$
\end{enumerate}
\end{myproof}
\end{prob}

% [Original: líneas 429-446 | solución completa | sin figura]
\begin{prob}
Sea $f(x,y)=\dfrac{x^4y}{x^8+y^2}.$
\begin{enumerate}[(a)]
    \item Determine los puntos donde la función es discontinua.
    \item Exprese si es posible remover las discontinuidades encontradas en el primer punto. Si lo es, determine los valores que hacen que la función sea continua. 
\end{enumerate}
\begin{myproof}
\begin{enumerate}[(a)]
\item Observe que la función solo se indetermina cuando $x^8+y^2=0,$ lo cual ocurre cuando $(x,y)=(0,0).$ De esta manera, al calcular el límite:
 \begin{align*}
\displaystyle\lim_{(x,y)\to (0,0)} \dfrac{x^4y}{x^8+y^2} {{=}\atop {\begin{array}{c}
\uparrow\\ \text{ si $y=mx^4$}\end{array}} } \displaystyle\lim_{x\to 0}\dfrac{x^4mx^4}{x^8+(mx^4)^2}=\displaystyle\lim_{x\to 0}\dfrac{x^8(m)}{x^8(1+m^2)}=\displaystyle\dfrac{m}{1+m^2}.
\end{align*}
Como este límite depende de $m,$ no existe.
\item No es posible remover las discontinuidades, pues el límite cuando la función tiende a $(0,0)$ no existe.
\end{enumerate}
\end{myproof}
\end{prob}

% [Original: líneas 685-702 | solución completa | sin figura]
\begin{prob}
Sea $f(x,y)=\dfrac{2xy^4}{x^5+6y^5}.$
\begin{enumerate}[(a)]
    \item Determine los puntos donde la función es discontinua.
    \item Exprese si es posible remover las discontinuidades encontradas en el primer punto. Si lo es, determine los valores que hacen que la función sea continua. 
\end{enumerate}
\begin{myproof}
\begin{enumerate}[(a)]
\item Observe que la función se indetermina cuando $x^5+6y^5=0.$ Esta ecuación no se reduce al origen: define la recta $x=-\sqrt[5]{6}\,y$ (pues $x^5=-6y^5$ implica $x=-6^{1/5}y$). En cualquier punto $(x_0,y_0)$ de esa recta con $y_0\neq 0$, el denominador es cero pero el numerador $2x_0y_0^4=-2\cdot 6^{1/5}y_0^5\neq 0$, produciendo una discontinuidad esencial (la función no está acotada). Adicionalmente, en el origen el límite no existe: tomando $y=mx$,
 \begin{align*}
\displaystyle\lim_{(x,y)\to (0,0)} \dfrac{2xy^4}{x^5+6y^5} {{=}\atop {\begin{array}{c}
\uparrow\\ \text{ si $y=mx$}\end{array}} } \displaystyle\lim_{x\to 0}\dfrac{2xm^4x^4}{x^5+6(mx)^5}=\displaystyle\lim_{x\to 0}\dfrac{x^5(2m^4)}{x^5(1+6m^5)}=\displaystyle\dfrac{2m^4}{1+6m^5},
\end{align*}
que depende de $m$ (para $1+6m^5\neq 0$), por lo que el límite en el origen tampoco existe.
\item No es posible remover ninguna de las discontinuidades: en los puntos de la recta $x=-\sqrt[5]{6}\,y$ con $y\neq 0$ la función diverge (discontinuidad esencial); en el origen el límite no existe.
\end{enumerate}
\end{myproof}
\end{prob}

% [Original: líneas 847-863 | solución completa | sin figura]
\begin{prob}
Sea $f(x,y)=\dfrac{8x^3y^2}{x^9+y^3}.$
\begin{enumerate}[(a)]
    \item Determine los puntos donde la función es discontinua.
    \item Exprese si es posible remover las discontinuidades encontradas en el primer punto. Si lo es, determine los valores que hacen que la función sea continua. 
\end{enumerate}
\begin{myproof}
\begin{enumerate}[(a)]
\item Observe que la función solo se indetermina cuando $x^9+y^3=0,$ lo cual ocurre cuando $y=-x^3.$ Suponga así que  $x=k$ entonces $y=-k^3.$ Así, al calcular el límite:
 \begin{align*}
\displaystyle\lim_{(x,y)\to (k,-k^3)} \dfrac{8x^3y^2}{x^9+y^3} = \displaystyle\lim_{x\to k}\dfrac{8k^3(-k^3)^2}{k^9+(-k^3)^3}&=\displaystyle\lim_{x\to k}\dfrac{8k^3(k^6)}{k^9+(-k^3)^3}\\&=\displaystyle\lim_{x\to k}\dfrac{8k^9}{k^9-k^9}=\displaystyle\dfrac{8k^9}{0}.
\end{align*}
El cual es un límite que no existe.
\item No es posible remover las discontinuidades sobre $y=-x^3$, pues el límite sobre los puntos de esta curva no existe.
\end{enumerate}
\end{myproof}
\end{prob}

% [Original: líneas 921-923 | sin verificar | sin figura]
\begin{prob}
Calcule $\displaystyle\lim_{(x,y)\to (0,0)}\left(\dfrac{x^4y^2+x^2y^4-5x^2-5y^2}{x^2+y^2}\right).$
\end{prob}

% [Original: líneas 945-947 | sin verificar | sin figura | 2 variante(s) de parámetro (se descartaron 1)]
\begin{prob}
Calcule $\displaystyle\lim_{(x,y)\to (5,9)}xy^2\left(\dfrac{x+4y}{x-y}\right).$
\end{prob}

% [Original: líneas 961-963 | sin verificar | sin figura]
\begin{prob}
Calcule $\displaystyle\lim_{(x,y)\to (2,1)}\ln\left(2x^2-3y^2\right).$
\end{prob}

% [Original: líneas 1020-1037 | solución completa | sin figura]
\begin{prob}
Sea $f(x,y)=\dfrac{x+y}{x-y}.$
\begin{enumerate}[(a)]
    \item Calcule $\displaystyle\lim_{(x,y)\to (0,0)}f(x)$ tomando la trayectoria $y=mx.$
    \item ¿Qué conclusión puede extraer?
\end{enumerate}
\begin{myproof}
Sea $f(x,y)=\dfrac{x+y}{x-y}.$ Observe que:
\begin{enumerate}[$a)$]
\item Tomando la trayectoria $y=mx$ el límite queda así:
\begin{align*}
\displaystyle\lim_{(x,y)\to (0,0)} \dfrac{x+y}{x-y} {{=}\atop {\begin{array}{c}
\uparrow\\ \text{ si $y=mx$}\end{array}} } \displaystyle\lim_{x\to 0}\dfrac{x+kx}{x-kx}=\displaystyle\lim_{x\to 0}\dfrac{x(1+k)}{x(1-k)}=\dfrac{1+k}{1-k}.
\end{align*}
\item El límite no existe, pues depende del valor de $k\in \mathbb{R}$.
\end{enumerate}
\end{myproof}
\end{prob}

% [Original: líneas 1063-1086 | solución completa | sin figura]
\begin{prob}
Sea $f(x,y)=\dfrac{x^4-y^2}{x^4+y^2}.$
\begin{enumerate}[(a)]
    \item Calcule $\displaystyle\lim_{(x,y)\to (0,0)}f(x)$ tomando la trayectoria $y=mx^2.$
    \item Calcule $\displaystyle\lim_{(x,y)\to (0,0)}f(x)$ tomando la trayectoria $y=0.$
    \item ¿Qué conclusión puede extraer?
\end{enumerate}
\begin{myproof}
Sea $f(x,y)=\dfrac{x^4-y^2}{x^4+y^2}.$ Observe que:
\begin{enumerate}[$a)$]
\item Tomando la trayectoria $y=mx^2$ el límite queda así:
\begin{align*}
\displaystyle\lim_{(x,y)\to (0,0)} \dfrac{x^4-y^2}{x^4+y^2} &{{=}\atop {\begin{array}{c}
\uparrow\\ \text{ si $y=mx^2$}\end{array}} } \displaystyle\lim_{x\to 0}\dfrac{x^4-m^2x^4}{x^4+m^2x^4}=\displaystyle\lim_{x\to 0}\dfrac{x^4(1-m^2)}{x^4(1+m^2)}\\&=\lim_{x\to 0}\dfrac{1-m^2}{1+m^2}=\dfrac{1-m^2}{1+m^2}.
\end{align*}
\item Tomando la trayectoria $y=0$ el límite queda así:
\begin{align*}
\displaystyle\lim_{(x,y)\to (0,0)} \dfrac{x^4-y^2}{x^4+y^2} {{=}\atop {\begin{array}{c}
\uparrow\\ \text{ si $y=0$}\end{array}} } \displaystyle\lim_{x\to 0}\dfrac{x^4}{x^4}=1.
\end{align*}
\item El límite no existe, pues dos trayectorias determinan un valor distinto.
\end{enumerate}
\end{myproof}
\end{prob}

% [Original: líneas 1111-1121 | solución completa | sin figura]
\begin{prob}
Sea $f(x,y)=x\cos\left(\dfrac{1}{y}\right).$ Sobre esta función se conoce que $-1\leq \cos\left(\dfrac{1}{y}\right)\leq 1$ Calcule $\displaystyle\lim_{(x,y)\to (0,0)}f(x,y).$
\begin{myproof}
Sea $f(x,y)=x\cos\left(\dfrac{1}{y}\right).$ Observe que:
\begin{enumerate}[$a)$]
\item $\displaystyle\lim_{(x,y)\to (0,0)} x = 0.$
\item $\displaystyle\lim_{(x,y)\to (0,0)} -x = 0.$
\item Por medio del teorema del emparedado el límite de la función $f(x,y)=x\cos\left(\dfrac{1}{y}\right)$  existe y es igual a $0.$
\end{enumerate}
\end{myproof}
\end{prob}

% [Original: líneas 1147-1154 | sin verificar | sin figura]
\begin{prob}
Dada la función $f(x,y)=\dfrac{x^3+y^3}{x^2+y^2}$ determine:
\begin{enumerate}[(a)]
\item El dominio de la función.
\item De acuerdo al resultado obtenido en el anterior ítem, ¿es posible redefinir la función $f(x,y)$ para que sea continua en todo $\mathbb{R}^2$?
\textit{Sugerencia:} Considere hacer una división de polinomios.
\end{enumerate}
\end{prob}

% [Original: líneas 1248-1250 | sin verificar | sin figura]
\begin{prob}
Determine si $\displaystyle\lim_{(x,y)\to (0,0)} \dfrac{Axy}{x^2+y^2}$ existe. Justifique adecuadamente su respuesta.
\end{prob}

% [Original: líneas 1380-1382 | sin verificar | sin figura]
\begin{prob}
Calcule $\displaystyle\lim_{(x,y)\to (0,0)}\left(\dfrac{x^3y+xy^3-3x^2-3y^2}{x^2+y^2}\right).$
\end{prob}

% [Original: líneas 1420-1422 | sin verificar | sin figura]
\begin{prob}
Calcule $\displaystyle\lim_{(x,y)\to (1,1)}\ln\left(2x^2-y^2\right).$
\end{prob}

